%=============================================================================
% Section 1: Introduction
%=============================================================================
\section{Introduction}

\subsection{The Integration Challenge}

AI applications increasingly depend on heterogeneous capabilities provided by model vendors, external services, and local devices. Providers such as OpenAI~\cite{openai2023}, Anthropic~\cite{anthropic2024}, Google, and DeepSeek expose powerful models through APIs, but direct integration introduces tight coupling and escalating system complexity.

Most current AI frameworks address this challenge through application-level orchestration (e.g., LangChain~\cite{langchain2024}). Although this approach improves development velocity, it frequently combines orchestration policy and execution mechanics in one runtime, thereby increasing lock-in and reducing component reusability.

\subsection{The Missing Layer}

The ecosystem still lacks a minimal execution layer dedicated to capability invocation. Such a layer should:
\begin{enumerate}[label=\arabic*.]
    \item Execute capabilities defined by standardized protocols
    \item Translate protocol requests into provider-specific API calls
    \item Normalize results into standardized responses
    \item Handle execution errors without implementing application-level retry strategies
\end{enumerate}

This layer excludes higher-level concerns such as planning, workflow orchestration, and application logic, thereby preserving a lightweight and reusable execution boundary.

\subsection{Contributions}

The primary contributions are as follows:
\begin{enumerate}[label=\arabic*.]
    \item \textbf{Formal Architectural Model:} A three-layer logical model $S = (A, P, E)$ is specified to separate cognition, orchestration, and execution with explicit layer responsibilities.
    \item \textbf{Minimal Runtime Constraint:} Minimality is formalized as a verifiable design constraint, with explicit runtime responsibilities and non-responsibilities for deterministic capability execution.
    \item \textbf{Unified Capability Semantics:} Capability is formalized as $C = (I, O, \Sigma, Q, P)$, enabling technical invocation and economic observability to share a single schema-level object.
    \item \textbf{Cross-Layer Integration:} The three-layer logical model is mapped to the five-layer AI-Native implementation stack, clarifying how execution-layer invariants are preserved in full-system deployments.
\end{enumerate}

The execution layer is then positioned within the broader five-layer AI-Native Architecture developed in later sections.

%-----------------------------------------------------------------------------
% Figure 1: Layered Architecture
%-----------------------------------------------------------------------------
\begin{figure}[htbp]
\centering
\begin{tikzpicture}[
    layerbox/.style={
        draw,
        rectangle,
        minimum width=10cm,
        minimum height=0.9cm,
        align=center,
        font=\small
    },
    layerbox tall/.style={
        draw,
        rectangle,
        minimum width=10cm,
        minimum height=2.2cm,
        align=center,
        font=\small
    },
    annot/.style={
        font=\footnotesize,
        text=gray!70!black
    },
    arrow/.style={
        -{Stealth[length=3mm, width=2mm]},
        thick,
        gray!60!black
    }
]

% Agent Intelligence Layer
\node[layerbox, fill=blue!8] (agent) at (0, 6) {Agent Intelligence Layer};
\node[annot] at (0, 5.35) {(Reasoning, Planning, Cognition)};

\draw[arrow] (0, 4.9) -- (0, 4.4);

% Contact Layer
\node[layerbox, fill=green!8] (contact) at (0, 3.8) {Contact Layer};
\node[annot] at (0, 3.15) {(Routing, Provider Selection, Workflow Control)};

\draw[arrow] (0, 2.7) -- (0, 2.2);

% AI Execution Layer (detailed box)
\node[layerbox tall, fill=orange!12] (exec) at (0, 0) {};
\node[font=\small\bfseries] at (0, 0.85) {AI Execution Layer (Runtime)};
\node[font=\footnotesize, align=left, text=gray!80!black] at (-0.5, 0.0) {
    \textbullet~Protocol $\rightarrow$ API Translation\\
    \textbullet~Result Normalization\\
    \textbullet~Connection Management
};
\node[font=\footnotesize, align=left, text=gray!80!black] at (-0.5, -0.85) {
    \textbullet~Error Propagation\\
    \textbullet~Context Forwarding
};

\draw[arrow] (0, -1.3) -- (0, -1.8);

% Protocol Boundary line
\node[font=\footnotesize\itshape, text=red!60!black] at (0, -2.1) {--- AI Protocol Boundary ---};

\draw[arrow] (0, -2.4) -- (0, -2.9);

% Capability Providers
\node[layerbox, fill=purple!8] (providers) at (0, -3.5) {Capability Providers / Tools / Devices};
\node[annot] at (0, -4.15) {(LLMs, APIs, Databases, External Services, IoT)};

\end{tikzpicture}
\caption{Layered architecture of the AI Execution Layer within AI-native systems, illustrating the separation between cognition, the Contact Layer, and capability execution.}
\label{fig:architecture}
\end{figure}
